\section{Example Question Metadata}
\label{app:question-metadata}

The original BEETLE and SciEntsBank datasets provide only a question and a reference answer as scoring context. We therefore augment each question with a more detailed question context and a rubric for each score level. \autoref{tab:question-metadata} presents one complete BEETLE example used to condition answer generation.

\begin{table*}[t]
    \centering
    \small
    \renewcommand{\arraystretch}{1.15}
    \resizebox{\linewidth}{!}{%
    \begin{tabular}{@{}p{0.16\textwidth}p{\dimexpr0.84\textwidth-2\tabcolsep\relax}@{}}
    \hline
    \textbf{Component} & \textbf{Content} \\
    \hline
    \raggedright\textbf{Question} &
    Finally, consider finding a burned out light bulb in a long string of lights. One approach would be to place one lead of the multimeter on a battery terminal and then move the other lead systematically, one light bulb terminal at a time. How would you know when you have found the burned out bulb? \\
    \hline
    \raggedright\textbf{Reference answer} &
    When there is a change in voltage you have found the damaged bulb \\
    \hline
    \raggedright\textbf{Question context} &
    Task: test a series circuit of light bulbs by holding one multimeter lead on a battery terminal and touching the other lead to each bulb terminal in turn. Relevant entities: battery, bulb terminals, multimeter voltage reading. Reasoning: a burned (open) bulb alters the circuit so the measured voltage suddenly changes at that bulb; that change identifies the fault. \\
    \hline
    \raggedright\textbf{Incorrect rubric} &
    Incorrect: Responses that state an absolute number or present no diagnostic criterion, report unrelated measurements, or give no indication of using a voltage change to locate the fault. Answers that assert a single fixed reading without connecting it to a change at a particular bulb, or that are irrelevant, earn zero. \\
    \hline
    \raggedright\textbf{Partially correct rubric} &
    Partially Correct: Answers that mention a voltage change but fail to explain how that change identifies the burned bulb earn partial credit. Credit for noting that the multimeter reading will differ at the faulty location, but missing clarification that the change occurs when moving from a working bulb terminal to the burned bulb terminal is partial. \\
    \hline
    \raggedright\textbf{Correct rubric} &
    Correct: Complete answers state that you detect the burned bulb by observing a change in the multimeter voltage as you move the probe along the string; the damaged bulb is located where the reading shifts, indicating a discontinuity or altered potential across that socket. Equivalent clear descriptions of finding the bulb by this voltage change earn full credit. \\
    \hline
    \end{tabular}%
    }
    \caption{Question, reference answer, question context, and three-level rubric for a BEETLE example.}
    \label{tab:question-metadata}
\end{table*}

\section{Example Synthesis Prompt}
\label{app:synthesis-prompt}

\autoref{tab:synthesis-prompt} presents a first-round synthesis prompt for the BEETLE question in \autoref{tab:question-metadata} at the partially correct level. The answers it lists are the generated seed examples for this question and level.

% TODO: check against the server code, since the local template says 5 answers and has no reference-answer field; confirm the "Real student answers" heading is also used with generated seed examples.
\begin{table*}[t]
    \centering
    \footnotesize
    \renewcommand{\arraystretch}{1.15}
    \begin{tabular}{@{}p{0.09\textwidth}p{\dimexpr0.91\textwidth-2\tabcolsep\relax}@{}}
    \hline
    \textbf{Role} & \textbf{Content} \\
    \hline
    \raggedright\textbf{System} &
    \raggedright\ttfamily You are an expert in educational assessment and student short-answer response writing.\newline\newline Generate realistic student answers for the requested score level. Match the language, style, brevity, imperfections, and conceptual quality of the provided examples. Return only the requested answers, one per line. \tabularnewline
    \hline
    \raggedright\textbf{User} &
    \raggedright\ttfamily Question:\newline
    Finally, consider finding a burned out light bulb in a long string of lights. One approach would be to place one lead of the multimeter on a battery terminal and then move the other lead systematically, one light bulb terminal at a time. How would you know when you have found the burned out bulb?\newline\newline
    Scoring context:\newline
    Task: test a series circuit of light bulbs by holding one multimeter lead on a battery terminal and touching the other lead to each bulb terminal in turn. Relevant entities: battery, bulb terminals, multimeter voltage reading. Reasoning: a burned (open) bulb alters the circuit so the measured voltage suddenly changes at that bulb; that change identifies the fault.\newline\newline
    Target score label:\newline
    Partially Correct\newline\newline
    Rubric for this question and target score label:\newline
    Partially Correct: Answers that mention a voltage change but fail to explain how that change identifies the burned bulb earn partial credit. Credit for noting that the multimeter reading will differ at the faulty location, but missing clarification that the change occurs when moving from a working bulb terminal to the burned bulb terminal is partial.\newline\newline
    Generated seed answers for this question and score label:\newline
    There is a change in voltage.\newline
    When the voltage reading changes.\newline
    The voltage will change.\newline
    A difference appears in voltage.\newline
    Where the voltage changes.\newline
    The meter shows changing voltage.\newline
    There is a voltage difference.\newline
    When the electrical state changes.\newline
    Voltage changes between the terminals.\newline
    The reading shows a voltage change.\newline\newline
    Task:\newline
    Write exactly 10 new synthesized student answers for the same question and target score label.
    \end{tabular}
    \caption{Synthesis prompt for a BEETLE question at the partially correct level.}
    \label{tab:synthesis-prompt}
\end{table*}

\section{Seed Examples for Unseen Questions}
\label{app:seed-examples}

Labeled answers to unseen questions are test data, and seeds written by an LLM from the question alone are not grounded in how students actually answer. We therefore use the answers to unseen questions without their labels and take level information only from the judge panel (\autoref{app:judge-scorer}), which is trained on the original training set. Judge inference and all LLM calls read label-stripped copies of the test splits that contain only the question, reference answer, and student answer.

\paragraph{Selection.}
For each answer $a$ and level $\ell$, the panel score is $s(a, \ell) = \frac{1}{J} \sum_{j=1}^{J} p_j(\ell \mid a)$, as in \autoref{eq:panel-score}. For each question and level, we select the 10 answers with the highest $s(a, \ell)$ among those whose argmax level is $\ell$, so each answer is selected for at most one level. Cells with fewer such answers keep only what is available.

\paragraph{Rewriting.}
For each question and target level, GPT-5.6-Luna receives the question, the reference answer, and the selected answers of all three levels grouped by predicted level, with the instruction that the grouping comes from a fallible classifier. It checks each selected answer against the reference answer, keeps it if it fits the target level, rewrites it otherwise, and infers new answers from the question and reference answer when fewer than 10 usable ones remain. Each call returns 10 short, informal answers, which become the seed examples for that question and level.

\paragraph{Question contexts and rubrics.}
The question contexts and rubrics for unseen questions (\autoref{app:question-metadata}) are generated by \textit{gpt-5-mini} from the question, the reference answer, and the selected answers grouped by predicted level, with the instruction that the rubrics remain consistent with the reference answer.

\section{Datasets and Evaluation Setup}
\label{app:datasets}

We use BEETLE and SciEntsBank~\citep{dzikovskaSemEval2013Task2013}. Each example consists of a question $q$, a reference answer $r$, a student answer $a$, and a human-assigned score level $\ell \in \{0, 1, 2\}$. The scorer predicts $\ell$ from $(q, r, a)$. Both benchmarks provide unseen questions (UQ) from the training domains, and SciEntsBank also provides questions from unseen domains (UD). These target questions are disjoint from the seen training questions. Their gold labels are reserved for evaluation, and their student answers are used for synthesis only without labels (\autoref{app:seed-examples}).

\begin{table}[t]
    \centering
    \small
    \begin{tabular*}{\linewidth}{@{\extracolsep{\fill}}lrr@{}}
    \hline
    \textbf{Split} & \textbf{Questions} & \textbf{Answers} \\
    \hline
    \multicolumn{3}{@{}l}{\textbf{BEETLE}} \\
    Train & 47 & 3,941 \\
    \sout{Test UA} & \sout{47} & \sout{439} \\
    Test UQ & 9 & 819 \\
    \hline
    \multicolumn{3}{@{}l}{\textbf{SciEntsBank}} \\
    Train & 135 & 4,969 \\
    \sout{Test UA} & \sout{135} & \sout{540} \\
    Test UQ & 15 & 733 \\
    Test UD & 46 & 4,562 \\
    \hline
    \end{tabular*}
    \caption{Number of questions and student answers in the original training and test splits. \sout{Test UA} is not used in the experiments.}
    \label{tab:dataset-statistics}
\end{table}

\section{Synthesizer Training Details}
\label{app:synthesizer-training}

\autoref{tab:synthesizer-hyperparameters} lists the hyperparameters of the three fine-tuned synthesizer variants. All variants are trained as LoRA adapters on Qwen2.5-7B-Instruct, Llama-3.1-8B-Instruct, and Mistral-7B-Instruct-v0.3, with identical settings for the three models and for BEETLE and SciEntsBank. We use the sigmoid preference loss and an NLL term on the chosen completion for DPO. The length tolerance is not part of the loss: rejected answers are matched to chosen answers one by one, and a pair is kept only if the token lengths of the two differ by at most 10\% of their mean.

\begin{table}[t]
    \centering
    \small
    \setlength{\tabcolsep}{4pt}
    \begin{tabular}{@{}lccc@{}}
    \hline
    & \textbf{SFT} & \textbf{SFT+DPO} & \textbf{DPO} \\
    \hline
    Starts from & base & SFT adapter & base \\
    $\pi_{\mathrm{ref}}$ & -- & SFT adapter & base \\
    LoRA rank / $\alpha_{\mathrm{LoRA}}$ & 32 / 64 & 32 / 64 & 64 / 64 \\
    LoRA dropout & 0.05 & 0.05 & 0.05 \\
    Learning rate & $1.5\times10^{-4}$ & $5\times10^{-6}$ & $1.5\times10^{-4}$ \\
    Steps & 800 & 400 & 1,200 \\
    Batch size & 32 & 16 & 16 \\
    Max length & 4,096 & 4,096 & 4,096 \\
    $\beta$ & -- & 0.1 & 0.1 \\
    $\alpha$ & -- & 1.0 & 1.0 \\
    Length tolerance & -- & 0.1 & 0.1 \\
    Checkpoint & last & best eval & best eval \\
    \hline
    \end{tabular}
    \caption{Hyperparameters of the fine-tuned synthesizers.}
    \label{tab:synthesizer-hyperparameters}
\end{table}

All runs apply LoRA to the query, key, value, and output projections and share the following settings: a cosine learning-rate schedule with 5\% warmup, the fused AdamW optimizer without weight decay, bf16 precision, gradient checkpointing, and random seed 13. Each training prompt shows no examples with probability 0.05 and 10 examples with probability 0.95, and each completion contains 10 answers.

\autoref{tab:generation-settings} lists the number of generation rounds and the sampling temperature of each synthesizer.

\begin{table}[t]
    \centering
    \small
    \begin{tabular}{@{}lccc@{}}
    \hline
    & \multicolumn{2}{c}{\textbf{Rounds}} & \\
    \cline{2-3}
    & \textbf{BEETLE} & \textbf{SciEntsBank} & \textbf{Temperature} \\
    \hline
    Prompting & 5 & 5 & 0.8 \\
    SFT & 19 & 10 & 1.0 \\
    SFT+DPO & 19 & 10 & 1.0 \\
    DPO & 19 & 10 & 1.0 \\
    \hline
    \end{tabular}
    \caption{Generation rounds and sampling temperature of each synthesizer. SciEntsBank UQ and UD use the same number of rounds.}
    \label{tab:generation-settings}
\end{table}

\section{Judge and Scorer Models}
\label{app:judge-scorer}

The judge panel consists of six sequence classifiers: BERT~\citep{devlinBERTPretrainingDeep2019}, RoBERTa~\citep{liuRoBERTaRobustlyOptimized2019}, ModernBERT~\citep{warnerSmarterBetterFaster2025}, DeBERTa~\citep{heDeBERTaV3ImprovingDeBERTa2023}, Llama-3.2-1B, and Llama-3.2-3B~\citep{grattafioriLlama3Herd2024}, with the Llama models used as encoders. Each judge reads the question, reference answer, and student answer and predicts the score level. The Llama judges are loaded with a sequence-classification head, which predicts the level from the hidden state of the last token, and are trained as LoRA adapters of rank 64 for 4 epochs with a learning rate of $10^{-4}$. The judges are trained only on the original training set, so they see neither synthetic answers nor answers to unseen questions. The scorer takes the same input, uses BERT, RoBERTa, ModernBERT, or DeBERTa as the encoder, and is trained with the same settings as the judge with the same encoder.
% TODO: add for BERT/RoBERTa/ModernBERT/DeBERTa: base or large, epochs, learning rate, seeds, checkpoint selection.

\autoref{tab:judge-f1-by-arm} reports how often the judges assign the synthetic answers of each synthesizer to the level they were generated for.

% Generated by analysis_outputs/judge_by_arm/plot_judge_by_arm.py.
\begin{table}[t]
\centering
\footnotesize
\setlength{\tabcolsep}{2pt}
\resizebox{\linewidth}{!}{%
\begin{tabular}{lcccccccc}
\toprule
 & \multicolumn{4}{c}{BEETLE} & \multicolumn{4}{c}{SciEntsBank} \\
\cmidrule(lr){2-5} \cmidrule(lr){6-9}
Synthesizer & Inc. & Part. & Corr. & Macro & Inc. & Part. & Corr. & Macro \\
\midrule
Prompting & 59.8 & 20.4 & 55.5 & 45.3 & 60.7 & 25.4 & 64.2 & 50.1 \\
SFT & 57.7 & 44.2 & 61.0 & 54.3 & 59.0 & 32.8 & 58.8 & 50.2 \\
SFT+DPO & 62.3 & 45.4 & 64.5 & 57.4 & 65.5 & 35.5 & 63.0 & 54.6 \\
DPO & \textbf{67.0} & \textbf{49.0} & \textbf{67.6} & \textbf{61.2} & \textbf{68.3} & \textbf{41.0} & \textbf{69.3} & \textbf{59.5} \\
\bottomrule
\end{tabular}
}%
\caption{F1 (\%) of the six judges on synthetic answers, against the level each was generated for. \textbf{Bold} marks the best per column.}
\label{tab:judge-f1-by-arm}
\end{table}

\section{Per-Model Scoring Results}
\label{app:per-model-results}

\autoref{tab:appendix-random-amount} reports the results for each scorer and synthesis method at every amount of synthetic data, and \autoref{tab:appendix-selection-amount} reports the same for each selection policy.

% Per-model appendix tables, generated by analysis_outputs/judge_uncertainty/build_appendix_tables.py
% Needs \usepackage{booktabs,multirow}.

\begin{table*}[t]
\centering
\footnotesize
\setlength{\tabcolsep}{2pt}
\begin{tabular}{llccccccccccccccc}
\toprule
 & & \multicolumn{5}{c}{BEETLE UQ} & \multicolumn{5}{c}{SciEntsBank UQ} & \multicolumn{5}{c}{SciEntsBank UD} \\
\cmidrule(lr){3-7}\cmidrule(lr){8-12}\cmidrule(lr){13-17}
Model & Synthesizer & 1$\times$ & 2$\times$ & 3$\times$ & 4$\times$ & 5$\times$ & 1$\times$ & 2$\times$ & 3$\times$ & 4$\times$ & 5$\times$ & 1$\times$ & 2$\times$ & 3$\times$ & 4$\times$ & 5$\times$ \\
\midrule
\multirow{5}{*}{ModernBERT} & w/o synthetic data & \multicolumn{5}{c}{61.8} & \multicolumn{5}{c}{60.1} & \multicolumn{5}{c}{54.6} \\
 & Prompting & \textcolor{red}{59.4} & \textcolor{red}{58.3} & \textcolor{red}{56.7} & \textcolor{red}{54.6} & \textcolor{red}{55.0} & \textcolor{red}{53.3} & \textcolor{red}{54.2} & \textcolor{red}{53.0} & \textcolor{red}{56.1} & \textcolor{red}{51.3} & \textcolor{red}{49.4} & \textcolor{red}{48.8} & \textcolor{red}{50.2} & \textcolor{red}{49.8} & \textcolor{red}{47.7} \\
 & SFT & 62.1 & \textcolor{red}{59.7} & 63.8 & \textcolor{red}{61.3} & \textcolor{red}{61.6} & \textcolor{red}{56.7} & \textcolor{red}{59.8} & 60.5 & \textcolor{red}{58.9} & \textcolor{red}{58.0} & 56.3 & 58.6 & 58.9 & 57.1 & 57.6 \\
 & SFT+DPO & \textcolor{red}{60.6} & \textcolor{red}{61.6} & 62.2 & 62.9 & 62.8 & \textcolor{red}{57.5} & \textcolor{red}{57.9} & 60.6 & 60.1 & 60.6 & 56.7 & 60.4 & 60.5 & 60.1 & 60.1 \\
 & DPO & 62.6 & 63.7 & 64.4 & 63.5 & 64.6 & 62.6 & 61.6 & 66.0 & 65.7 & 66.4 & 60.2 & 60.5 & 61.3 & 61.9 & 62.7 \\
\midrule
\multirow{5}{*}{BERT} & w/o synthetic data & \multicolumn{5}{c}{56.0} & \multicolumn{5}{c}{43.9} & \multicolumn{5}{c}{47.6} \\
 & Prompting & 60.0 & 56.5 & \textcolor{red}{55.3} & \textcolor{red}{51.8} & \textcolor{red}{53.0} & 49.3 & 48.0 & 49.4 & 48.7 & 47.1 & \textcolor{red}{44.1} & \textcolor{red}{44.3} & \textcolor{red}{43.8} & \textcolor{red}{45.0} & \textcolor{red}{44.9} \\
 & SFT & 60.1 & 61.4 & 61.2 & 60.5 & 60.6 & 52.2 & 52.4 & 54.5 & 54.0 & 54.9 & 52.9 & 54.0 & 54.9 & 55.2 & 56.1 \\
 & SFT+DPO & 59.7 & 60.0 & 61.7 & 60.8 & 61.9 & 52.3 & 53.6 & 56.1 & 54.5 & 55.0 & 54.0 & 56.1 & 56.8 & 57.0 & 57.8 \\
 & DPO & 61.5 & 61.9 & 61.9 & 61.8 & 61.9 & 56.3 & 57.9 & 59.0 & 59.8 & 60.9 & 55.4 & 57.2 & 58.0 & 58.4 & 57.6 \\
\midrule
\multirow{5}{*}{RoBERTa} & w/o synthetic data & \multicolumn{5}{c}{61.9} & \multicolumn{5}{c}{58.9} & \multicolumn{5}{c}{54.0} \\
 & Prompting & 62.3 & \textcolor{red}{58.4} & \textcolor{red}{57.9} & \textcolor{red}{57.7} & \textcolor{red}{56.2} & \textcolor{red}{55.2} & \textcolor{red}{55.5} & \textcolor{red}{54.9} & \textcolor{red}{53.5} & \textcolor{red}{54.3} & 54.7 & \textcolor{red}{51.8} & \textcolor{red}{50.5} & \textcolor{red}{50.2} & \textcolor{red}{49.7} \\
 & SFT & 62.7 & 63.5 & 63.4 & 63.8 & 65.0 & \textcolor{red}{57.5} & \textcolor{red}{58.5} & 60.6 & 60.8 & 59.1 & 58.8 & 59.3 & 59.2 & 59.5 & 59.4 \\
 & SFT+DPO & 62.5 & 63.0 & 63.4 & 63.2 & 63.6 & 59.5 & 60.0 & 61.4 & 59.8 & 60.9 & 59.7 & 60.4 & 61.2 & 60.4 & 60.9 \\
 & DPO & 65.4 & 65.4 & 64.4 & 65.9 & 64.4 & 65.9 & 64.5 & 65.8 & 65.3 & 63.9 & 60.0 & 60.6 & 61.5 & 61.2 & 60.2 \\
\midrule
\multirow{5}{*}{DeBERTa} & w/o synthetic data & \multicolumn{5}{c}{65.1} & \multicolumn{5}{c}{63.0} & \multicolumn{5}{c}{58.8} \\
 & Prompting & \textcolor{red}{61.7} & \textcolor{red}{58.6} & \textcolor{red}{56.6} & \textcolor{red}{55.4} & \textcolor{red}{54.7} & \textcolor{red}{56.9} & \textcolor{red}{57.2} & \textcolor{red}{56.2} & \textcolor{red}{55.7} & \textcolor{red}{56.0} & \textcolor{red}{54.2} & \textcolor{red}{52.2} & \textcolor{red}{53.6} & \textcolor{red}{52.4} & \textcolor{red}{51.3} \\
 & SFT & \textcolor{red}{63.3} & \textcolor{red}{63.7} & \textcolor{red}{64.4} & \textcolor{red}{63.7} & \textcolor{red}{62.9} & \textcolor{red}{61.4} & \textcolor{red}{59.6} & \textcolor{red}{62.2} & \textcolor{red}{61.8} & \textcolor{red}{61.4} & 59.5 & 60.6 & 59.5 & 59.3 & 58.8 \\
 & SFT+DPO & \textcolor{red}{63.6} & \textcolor{red}{63.4} & \textcolor{red}{63.0} & \textcolor{red}{62.7} & \textcolor{red}{64.3} & \textcolor{red}{61.9} & \textcolor{red}{60.6} & \textcolor{red}{62.6} & \textcolor{red}{61.6} & \textcolor{red}{61.8} & 61.1 & 61.3 & 62.2 & 61.1 & 61.0 \\
 & DPO & 66.4 & \textcolor{red}{64.7} & 65.9 & 65.8 & 66.5 & 65.6 & 66.6 & 66.4 & 67.2 & 66.8 & 62.7 & 62.4 & 63.0 & 62.7 & 62.6 \\
\bottomrule
\end{tabular}
\caption{Macro-F1 ($\times$100) of each scorer trained on real data plus a random sample of synthetic answers from each synthesizer, by amount of synthetic data (1$\times$--5$\times$). Synthetic cells are means over 5 sampling draws; the score without synthetic data spans each test set; red marks cells below it.}
\label{tab:appendix-random-amount}
\end{table*}

\begin{table*}[t]
\centering
\footnotesize
\setlength{\tabcolsep}{2pt}
\begin{tabular}{llccccccccccccccc}
\toprule
 & & \multicolumn{5}{c}{BEETLE UQ} & \multicolumn{5}{c}{SciEntsBank UQ} & \multicolumn{5}{c}{SciEntsBank UD} \\
\cmidrule(lr){3-7}\cmidrule(lr){8-12}\cmidrule(lr){13-17}
Model & Selection & 1$\times$ & 2$\times$ & 3$\times$ & 4$\times$ & 5$\times$ & 1$\times$ & 2$\times$ & 3$\times$ & 4$\times$ & 5$\times$ & 1$\times$ & 2$\times$ & 3$\times$ & 4$\times$ & 5$\times$ \\
\midrule
\multirow{9}{*}{ModernBERT} & w/o synthetic data & \multicolumn{5}{c}{61.8} & \multicolumn{5}{c}{60.1} & \multicolumn{5}{c}{54.6} \\
 & Random & \textcolor{red}{61.6} & 62.6 & 63.3 & 63.2 & 63.7 & 60.1 & \textcolor{red}{59.7} & 63.3 & 62.9 & 63.5 & 58.5 & 60.4 & 60.9 & 61.0 & 61.4 \\
 & Highest-Rank & 63.2 & -- & -- & -- & -- & 62.0 & -- & -- & -- & -- & 56.9 & -- & -- & -- & -- \\
 & Lowest-Rank & \textcolor{red}{44.8} & -- & -- & -- & -- & \textcolor{red}{36.8} & -- & -- & -- & -- & \textcolor{red}{35.9} & -- & -- & -- & -- \\
 & Hard Selection & 65.2 & 65.0 & 65.1 & 65.1 & 65.1 & 65.5 & 65.3 & 65.5 & 65.4 & 65.3 & 61.6 & 62.2 & 62.4 & 62.8 & 63.2 \\
 & $(0, 30)$ & 64.0 & 65.3 & 66.1 & 65.8 & 66.6 & 63.4 & 62.8 & 64.5 & 64.8 & 64.5 & 62.0 & 62.6 & 63.5 & 63.8 & 63.8 \\
 & $(0, 50)$ & 64.7 & 65.5 & 65.5 & 65.9 & 65.5 & 63.9 & 65.1 & 65.0 & 64.7 & 64.9 & 62.2 & 62.8 & 63.1 & 63.1 & 63.5 \\
 & $(10, 50)$ & 64.6 & 65.5 & 65.4 & 65.9 & 66.0 & 64.0 & 64.2 & 65.0 & 65.2 & 64.9 & 62.2 & 62.5 & 63.6 & 63.4 & 63.6 \\
 & $(20, 20)$ & 64.5 & 66.2 & 65.9 & 66.4 & 66.4 & 61.9 & 62.7 & 63.8 & 62.5 & 64.0 & 60.6 & 62.3 & 63.9 & 62.0 & 63.5 \\
\midrule
\multirow{9}{*}{BERT} & w/o synthetic data & \multicolumn{5}{c}{56.0} & \multicolumn{5}{c}{43.9} & \multicolumn{5}{c}{47.6} \\
 & Random & 60.6 & 60.9 & 61.8 & 61.3 & 61.9 & 54.3 & 55.7 & 57.5 & 57.2 & 57.9 & 54.7 & 56.7 & 57.4 & 57.7 & 57.7 \\
 & Highest-Rank & 60.9 & -- & -- & -- & -- & 55.3 & -- & -- & -- & -- & 52.1 & -- & -- & -- & -- \\
 & Lowest-Rank & \textcolor{red}{41.6} & -- & -- & -- & -- & \textcolor{red}{38.1} & -- & -- & -- & -- & \textcolor{red}{35.1} & -- & -- & -- & -- \\
 & Hard Selection & 64.5 & 64.4 & 64.6 & 64.5 & 64.2 & 56.2 & 59.2 & 60.4 & 61.0 & 61.0 & 56.6 & 59.5 & 60.1 & 61.0 & 61.3 \\
 & $(0, 30)$ & 63.4 & 64.7 & 65.4 & 65.2 & 65.2 & 56.0 & 57.7 & 59.0 & 60.3 & 60.7 & 57.2 & 60.6 & 61.3 & 61.2 & 62.3 \\
 & $(0, 50)$ & 63.8 & 64.6 & 64.6 & 64.9 & 64.6 & 56.7 & 58.5 & 60.3 & 60.3 & 61.5 & 57.8 & 60.0 & 61.3 & 61.6 & 61.9 \\
 & $(10, 50)$ & 64.1 & 64.8 & 65.2 & 65.2 & 65.2 & 56.5 & 59.0 & 60.3 & 60.5 & 61.1 & 58.4 & 60.9 & 61.4 & 61.7 & 61.9 \\
 & $(20, 20)$ & 62.5 & 63.3 & 65.4 & 64.3 & 65.7 & 55.4 & 56.8 & 57.7 & 58.5 & 59.7 & 56.2 & 59.1 & 60.0 & 60.2 & 60.7 \\
\midrule
\multirow{9}{*}{RoBERTa} & w/o synthetic data & \multicolumn{5}{c}{61.9} & \multicolumn{5}{c}{58.9} & \multicolumn{5}{c}{54.0} \\
 & Random & 64.0 & 64.2 & 63.9 & 64.6 & 64.0 & 62.7 & 62.3 & 63.6 & 62.5 & 62.4 & 59.9 & 60.5 & 61.3 & 60.8 & 60.5 \\
 & Highest-Rank & 64.3 & -- & -- & -- & -- & 61.2 & -- & -- & -- & -- & 55.8 & -- & -- & -- & -- \\
 & Lowest-Rank & \textcolor{red}{44.4} & -- & -- & -- & -- & \textcolor{red}{42.5} & -- & -- & -- & -- & \textcolor{red}{38.2} & -- & -- & -- & -- \\
 & Hard Selection & 65.1 & 64.9 & 64.7 & 64.8 & 64.7 & 64.2 & 64.6 & 65.9 & 66.0 & 65.8 & 60.9 & 61.8 & 62.2 & 62.4 & 62.8 \\
 & $(0, 30)$ & 65.4 & 65.9 & 66.4 & 66.9 & 67.0 & 63.9 & 63.9 & 65.4 & 65.0 & 65.8 & 62.1 & 63.3 & 63.5 & 63.6 & 63.8 \\
 & $(0, 50)$ & 65.1 & 65.1 & 66.2 & 66.2 & 66.1 & 63.5 & 64.4 & 65.1 & 65.2 & 65.5 & 61.8 & 62.9 & 63.1 & 63.0 & 62.9 \\
 & $(10, 50)$ & 65.0 & 65.6 & 66.6 & 66.5 & 65.8 & 63.2 & 64.5 & 65.1 & 65.1 & 65.8 & 61.2 & 62.7 & 62.8 & 63.1 & 63.1 \\
 & $(20, 20)$ & 65.0 & 66.9 & 67.5 & 66.6 & 66.1 & 62.5 & 63.8 & 64.4 & 64.2 & 64.3 & 61.7 & 63.1 & 62.9 & 62.9 & 62.5 \\
\midrule
\multirow{9}{*}{DeBERTa} & w/o synthetic data & \multicolumn{5}{c}{65.1} & \multicolumn{5}{c}{63.0} & \multicolumn{5}{c}{58.8} \\
 & Random & \textcolor{red}{65.0} & \textcolor{red}{64.0} & \textcolor{red}{64.4} & \textcolor{red}{64.3} & 65.4 & 63.8 & 63.6 & 64.5 & 64.4 & 64.3 & 61.9 & 61.9 & 62.6 & 61.9 & 61.8 \\
 & Highest-Rank & \textcolor{red}{63.8} & -- & -- & -- & -- & 64.2 & -- & -- & -- & -- & 60.1 & -- & -- & -- & -- \\
 & Lowest-Rank & \textcolor{red}{43.6} & -- & -- & -- & -- & \textcolor{red}{41.1} & -- & -- & -- & -- & \textcolor{red}{37.7} & -- & -- & -- & -- \\
 & Hard Selection & 66.3 & 65.3 & 66.2 & 65.6 & 66.2 & 65.6 & 66.2 & 66.6 & 66.3 & 66.4 & 62.5 & 64.4 & 63.9 & 63.7 & 63.9 \\
 & $(0, 30)$ & 66.1 & 67.2 & 67.8 & 67.8 & 67.5 & 63.8 & 65.0 & 66.5 & 65.7 & 66.7 & 63.7 & 65.2 & 65.0 & 64.8 & 64.4 \\
 & $(0, 50)$ & 65.3 & 67.0 & 66.5 & 66.8 & 66.4 & 65.4 & 66.0 & 65.5 & 66.8 & 66.3 & 64.0 & 64.5 & 64.5 & 64.5 & 64.6 \\
 & $(10, 50)$ & 66.5 & 66.6 & 67.0 & 67.1 & 66.9 & 65.3 & 65.0 & 66.5 & 66.2 & 66.8 & 64.1 & 64.4 & 64.3 & 64.4 & 64.5 \\
 & $(20, 20)$ & 65.9 & 67.0 & 67.5 & 68.0 & 67.4 & 64.1 & 63.9 & 65.3 & 64.5 & 65.4 & 63.7 & 63.6 & 64.7 & 64.0 & 63.6 \\
\bottomrule
\end{tabular}
\caption{Macro-F1 ($\times$100) of each scorer with each selection policy, by amount of synthetic data (1$\times$--5$\times$). Cells are means over the SFT+DPO and DPO synthesizers and 5 sampling draws; Highest-Rank and Lowest-Rank were run at 1$\times$ only. The score without synthetic data spans each test set; red marks cells below it.}
\label{tab:appendix-selection-amount}
\end{table*}


\autoref{tab:appendix-random-amount-level-0}, \autoref{tab:appendix-random-amount-level-1}, and \autoref{tab:appendix-random-amount-level-2} report the corresponding per-level F1 scores for incorrect, partially correct, and correct answers, respectively. \autoref{tab:appendix-selection-amount-level-0}, \autoref{tab:appendix-selection-amount-level-1}, and \autoref{tab:appendix-selection-amount-level-2} give the same breakdown for each selection policy.

\begin{table*}[t]
\centering
\footnotesize
\setlength{\tabcolsep}{2pt}
\begin{tabular}{llccccccccccccccc}
\toprule
 & & \multicolumn{5}{c}{BEETLE UQ} & \multicolumn{5}{c}{SciEntsBank UQ} & \multicolumn{5}{c}{SciEntsBank UD} \\
\cmidrule(lr){3-7}\cmidrule(lr){8-12}\cmidrule(lr){13-17}
Model & Synthesizer & 1$\times$ & 2$\times$ & 3$\times$ & 4$\times$ & 5$\times$ & 1$\times$ & 2$\times$ & 3$\times$ & 4$\times$ & 5$\times$ & 1$\times$ & 2$\times$ & 3$\times$ & 4$\times$ & 5$\times$ \\
\midrule
\multirow{5}{*}{ModernBERT} & w/o synthetic data & \multicolumn{5}{c}{68.6} & \multicolumn{5}{c}{64.6} & \multicolumn{5}{c}{64.2} \\
 & Prompting & 70.7 & 71.0 & 70.4 & 70.2 & 71.8 & 66.8 & 67.2 & 66.9 & 69.7 & 65.4 & 67.4 & 67.3 & 67.7 & 68.7 & 66.1 \\
 & SFT & 70.3 & \textcolor{red}{68.3} & 71.8 & 70.8 & 70.9 & 65.3 & 66.8 & 68.6 & 68.0 & 67.0 & 67.1 & 69.5 & 69.6 & 69.0 & 69.9 \\
 & SFT+DPO & 69.8 & 70.8 & 70.7 & 71.1 & 71.3 & 66.4 & 68.7 & 68.7 & 69.5 & 69.0 & 67.1 & 70.5 & 70.5 & 70.9 & 70.6 \\
 & DPO & 71.3 & 73.6 & 74.5 & 73.3 & 74.7 & 67.0 & 67.7 & 68.8 & 72.1 & 71.5 & 68.8 & 68.8 & 69.3 & 70.7 & 70.8 \\
\midrule
\multirow{5}{*}{BERT} & w/o synthetic data & \multicolumn{5}{c}{63.5} & \multicolumn{5}{c}{49.9} & \multicolumn{5}{c}{57.2} \\
 & Prompting & 67.0 & 66.5 & 65.6 & 65.8 & 66.1 & 54.5 & 57.8 & 57.9 & 57.7 & 58.4 & 60.1 & 61.1 & 60.3 & 60.8 & 60.8 \\
 & SFT & 66.4 & 67.0 & 68.4 & 67.6 & 66.1 & 55.9 & 57.5 & 59.9 & 59.8 & 60.9 & 61.1 & 62.4 & 62.7 & 64.6 & 64.8 \\
 & SFT+DPO & 66.3 & 66.4 & 68.3 & 65.9 & 67.7 & 55.5 & 59.1 & 62.3 & 61.2 & 61.4 & 62.2 & 64.5 & 65.6 & 67.3 & 66.2 \\
 & DPO & 67.4 & 69.1 & 69.2 & 70.0 & 70.5 & 53.8 & 58.3 & 58.8 & 62.0 & 63.6 & 59.7 & 62.2 & 63.8 & 64.4 & 63.3 \\
\midrule
\multirow{5}{*}{RoBERTa} & w/o synthetic data & \multicolumn{5}{c}{71.1} & \multicolumn{5}{c}{66.5} & \multicolumn{5}{c}{66.9} \\
 & Prompting & \textcolor{red}{70.7} & \textcolor{red}{69.6} & \textcolor{red}{69.8} & \textcolor{red}{70.8} & \textcolor{red}{69.0} & 66.9 & 67.6 & \textcolor{red}{65.7} & \textcolor{red}{65.5} & \textcolor{red}{65.5} & 68.6 & 67.5 & \textcolor{red}{66.4} & \textcolor{red}{66.8} & \textcolor{red}{66.8} \\
 & SFT & \textcolor{red}{70.4} & 71.4 & 71.9 & \textcolor{red}{70.9} & 73.5 & 66.8 & 66.8 & 68.6 & 68.2 & 67.8 & 68.2 & 68.6 & 69.1 & 69.7 & 69.9 \\
 & SFT+DPO & 71.2 & 71.3 & 71.9 & 71.8 & 72.1 & 67.3 & 68.4 & 69.5 & 68.5 & 69.6 & 69.2 & 70.2 & 71.5 & 71.6 & 71.0 \\
 & DPO & 73.9 & 75.5 & 74.7 & 75.9 & 74.4 & 68.4 & 69.5 & 68.2 & 69.6 & 67.5 & 68.1 & 69.3 & 69.2 & 69.4 & 68.9 \\
\midrule
\multirow{5}{*}{DeBERTa} & w/o synthetic data & \multicolumn{5}{c}{74.4} & \multicolumn{5}{c}{69.4} & \multicolumn{5}{c}{70.3} \\
 & Prompting & \textcolor{red}{71.4} & \textcolor{red}{71.5} & \textcolor{red}{70.9} & \textcolor{red}{70.4} & \textcolor{red}{70.2} & 69.6 & 70.6 & 69.7 & \textcolor{red}{68.9} & \textcolor{red}{68.9} & 72.2 & 71.6 & 71.5 & \textcolor{red}{70.3} & \textcolor{red}{69.3} \\
 & SFT & \textcolor{red}{71.0} & \textcolor{red}{72.3} & \textcolor{red}{73.0} & \textcolor{red}{71.7} & \textcolor{red}{71.2} & 69.5 & 70.1 & 71.9 & 71.6 & 71.3 & 70.9 & 72.3 & 71.1 & 71.9 & 70.6 \\
 & SFT+DPO & \textcolor{red}{72.1} & \textcolor{red}{73.2} & \textcolor{red}{71.2} & \textcolor{red}{69.8} & \textcolor{red}{71.3} & \textcolor{red}{69.4} & 71.5 & 72.5 & 70.2 & 71.8 & 71.6 & 72.7 & 73.2 & 72.5 & 72.6 \\
 & DPO & 75.2 & 75.4 & 74.9 & 75.6 & 76.3 & 69.9 & 71.5 & 71.1 & 72.7 & 71.9 & 70.5 & 70.7 & 71.9 & 71.8 & 71.5 \\
\bottomrule
\end{tabular}
\caption{F1 for incorrect answers ($\times$100) of each scorer trained on real data plus a random sample of synthetic answers from each synthesizer, by amount of synthetic data (1$\times$--5$\times$). Synthetic cells are means over five sampling draws; the per-level score without synthetic data spans each test set; \textcolor{red}{red} marks cells below it.}
\label{tab:appendix-random-amount-level-0}
\end{table*}

\begin{table*}[t]
\centering
\footnotesize
\setlength{\tabcolsep}{2pt}
\begin{tabular}{llccccccccccccccc}
\toprule
 & & \multicolumn{5}{c}{BEETLE UQ} & \multicolumn{5}{c}{SciEntsBank UQ} & \multicolumn{5}{c}{SciEntsBank UD} \\
\cmidrule(lr){3-7}\cmidrule(lr){8-12}\cmidrule(lr){13-17}
Model & Synthesizer & 1$\times$ & 2$\times$ & 3$\times$ & 4$\times$ & 5$\times$ & 1$\times$ & 2$\times$ & 3$\times$ & 4$\times$ & 5$\times$ & 1$\times$ & 2$\times$ & 3$\times$ & 4$\times$ & 5$\times$ \\
\midrule
\multirow{5}{*}{ModernBERT} & w/o synthetic data & \multicolumn{5}{c}{47.8} & \multicolumn{5}{c}{42.0} & \multicolumn{5}{c}{34.4} \\
 & Prompting & \textcolor{red}{47.2} & \textcolor{red}{46.8} & \textcolor{red}{45.6} & \textcolor{red}{44.7} & \textcolor{red}{44.4} & \textcolor{red}{39.4} & 44.8 & 45.6 & 46.9 & 44.1 & 40.3 & 42.2 & 44.5 & 45.0 & 43.4 \\
 & SFT & 48.5 & \textcolor{red}{46.3} & 50.4 & \textcolor{red}{46.5} & \textcolor{red}{47.6} & \textcolor{red}{41.4} & 43.5 & 45.7 & 43.2 & \textcolor{red}{41.9} & 42.6 & 44.0 & 45.1 & 41.8 & 41.0 \\
 & SFT+DPO & \textcolor{red}{45.9} & \textcolor{red}{46.4} & 48.9 & 49.2 & 49.2 & \textcolor{red}{41.8} & \textcolor{red}{38.6} & 45.9 & 44.6 & 45.1 & 44.4 & 46.4 & 46.3 & 45.3 & 45.4 \\
 & DPO & 49.6 & 50.0 & 50.1 & 49.8 & 51.1 & 48.3 & 44.8 & 52.8 & 51.4 & 51.8 & 45.5 & 44.2 & 44.5 & 45.3 & 46.3 \\
\midrule
\multirow{5}{*}{BERT} & w/o synthetic data & \multicolumn{5}{c}{37.8} & \multicolumn{5}{c}{16.0} & \multicolumn{5}{c}{24.9} \\
 & Prompting & 50.5 & 44.3 & 45.4 & 42.0 & 43.4 & 39.1 & 37.6 & 37.9 & 37.9 & 35.1 & 37.0 & 38.2 & 39.0 & 40.1 & 40.2 \\
 & SFT & 48.1 & 50.4 & 49.4 & 48.1 & 49.5 & 38.5 & 39.5 & 38.0 & 40.0 & 39.3 & 39.7 & 41.5 & 39.9 & 39.5 & 41.4 \\
 & SFT+DPO & 48.8 & 47.8 & 50.2 & 48.5 & 49.0 & 39.4 & 37.2 & 41.2 & 39.2 & 40.1 & 42.0 & 40.9 & 41.7 & 41.7 & 42.2 \\
 & DPO & 49.9 & 50.9 & 51.0 & 50.5 & 49.9 & 47.1 & 46.3 & 46.6 & 47.2 & 46.9 & 40.5 & 42.6 & 40.7 & 42.2 & 40.6 \\
\midrule
\multirow{5}{*}{RoBERTa} & w/o synthetic data & \multicolumn{5}{c}{44.8} & \multicolumn{5}{c}{36.7} & \multicolumn{5}{c}{32.1} \\
 & Prompting & 49.6 & 46.4 & 46.0 & 46.8 & 45.1 & 39.6 & 41.7 & 42.0 & 39.7 & 42.4 & 46.4 & 45.7 & 44.9 & 45.6 & 45.5 \\
 & SFT & 48.5 & 49.5 & 48.1 & 50.1 & 51.5 & 39.0 & 39.3 & 43.9 & 45.5 & 41.3 & 42.7 & 42.6 & 41.8 & 41.9 & 41.7 \\
 & SFT+DPO & 48.4 & 48.7 & 48.2 & 48.9 & 48.0 & 42.5 & 41.7 & 44.2 & 42.7 & 44.7 & 45.0 & 43.3 & 44.2 & 43.6 & 44.9 \\
 & DPO & 51.7 & 51.3 & 49.9 & 51.9 & 50.8 & 52.3 & 49.2 & 52.4 & 50.5 & 49.4 & 42.6 & 42.5 & 43.7 & 42.3 & 40.2 \\
\midrule
\multirow{5}{*}{DeBERTa} & w/o synthetic data & \multicolumn{5}{c}{45.6} & \multicolumn{5}{c}{45.1} & \multicolumn{5}{c}{38.1} \\
 & Prompting & 47.5 & 46.0 & \textcolor{red}{45.2} & \textcolor{red}{43.3} & \textcolor{red}{44.0} & \textcolor{red}{45.1} & 46.1 & 45.2 & \textcolor{red}{44.6} & 45.9 & 46.6 & 46.5 & 47.9 & 47.2 & 46.8 \\
 & SFT & 47.9 & 48.5 & 49.9 & 48.6 & 48.1 & 46.0 & \textcolor{red}{42.0} & 46.9 & 45.6 & 46.9 & 43.2 & 46.0 & 42.6 & 43.1 & 42.9 \\
 & SFT+DPO & 47.4 & 46.5 & 47.9 & 48.6 & 48.9 & 45.9 & \textcolor{red}{43.1} & 46.2 & 45.1 & 47.4 & 45.0 & 46.0 & 45.7 & 44.4 & 44.9 \\
 & DPO & 51.1 & 49.9 & 51.0 & 50.8 & 51.2 & 52.4 & 53.3 & 52.1 & 53.0 & 54.0 & 47.0 & 45.0 & 44.1 & 44.0 & 44.1 \\
\bottomrule
\end{tabular}
\caption{F1 for partially correct answers ($\times$100) of each scorer trained on real data plus a random sample of synthetic answers from each synthesizer, by amount of synthetic data (1$\times$--5$\times$). Synthetic cells are means over five sampling draws; the per-level score without synthetic data spans each test set; \textcolor{red}{red} marks cells below it.}
\label{tab:appendix-random-amount-level-1}
\end{table*}

\begin{table*}[t]
\centering
\footnotesize
\setlength{\tabcolsep}{2pt}
\begin{tabular}{llccccccccccccccc}
\toprule
 & & \multicolumn{5}{c}{BEETLE UQ} & \multicolumn{5}{c}{SciEntsBank UQ} & \multicolumn{5}{c}{SciEntsBank UD} \\
\cmidrule(lr){3-7}\cmidrule(lr){8-12}\cmidrule(lr){13-17}
Model & Synthesizer & 1$\times$ & 2$\times$ & 3$\times$ & 4$\times$ & 5$\times$ & 1$\times$ & 2$\times$ & 3$\times$ & 4$\times$ & 5$\times$ & 1$\times$ & 2$\times$ & 3$\times$ & 4$\times$ & 5$\times$ \\
\midrule
\multirow{5}{*}{ModernBERT} & w/o synthetic data & \multicolumn{5}{c}{69.0} & \multicolumn{5}{c}{73.8} & \multicolumn{5}{c}{65.3} \\
 & Prompting & \textcolor{red}{60.3} & \textcolor{red}{57.2} & \textcolor{red}{54.1} & \textcolor{red}{48.9} & \textcolor{red}{48.8} & \textcolor{red}{53.6} & \textcolor{red}{50.7} & \textcolor{red}{46.4} & \textcolor{red}{51.6} & \textcolor{red}{44.5} & \textcolor{red}{40.6} & \textcolor{red}{36.9} & \textcolor{red}{38.3} & \textcolor{red}{35.7} & \textcolor{red}{33.6} \\
 & SFT & \textcolor{red}{67.5} & \textcolor{red}{64.7} & 69.2 & \textcolor{red}{66.4} & \textcolor{red}{66.4} & \textcolor{red}{63.5} & \textcolor{red}{69.1} & \textcolor{red}{67.1} & \textcolor{red}{65.4} & \textcolor{red}{65.0} & \textcolor{red}{59.2} & \textcolor{red}{62.3} & \textcolor{red}{62.0} & \textcolor{red}{60.7} & \textcolor{red}{62.0} \\
 & SFT+DPO & \textcolor{red}{66.0} & \textcolor{red}{67.6} & \textcolor{red}{67.1} & \textcolor{red}{68.4} & \textcolor{red}{67.9} & \textcolor{red}{64.3} & \textcolor{red}{66.3} & \textcolor{red}{67.1} & \textcolor{red}{66.3} & \textcolor{red}{67.8} & \textcolor{red}{58.6} & \textcolor{red}{64.2} & \textcolor{red}{64.5} & \textcolor{red}{64.0} & \textcolor{red}{64.3} \\
 & DPO & \textcolor{red}{67.0} & \textcolor{red}{67.4} & \textcolor{red}{68.4} & \textcolor{red}{67.4} & \textcolor{red}{68.0} & \textcolor{red}{72.6} & \textcolor{red}{72.4} & 76.3 & \textcolor{red}{73.6} & 75.9 & 66.4 & 68.4 & 70.2 & 69.8 & 71.0 \\
\midrule
\multirow{5}{*}{BERT} & w/o synthetic data & \multicolumn{5}{c}{66.7} & \multicolumn{5}{c}{65.7} & \multicolumn{5}{c}{60.6} \\
 & Prompting & \textcolor{red}{62.5} & \textcolor{red}{58.7} & \textcolor{red}{54.8} & \textcolor{red}{47.6} & \textcolor{red}{49.6} & \textcolor{red}{54.2} & \textcolor{red}{48.8} & \textcolor{red}{52.3} & \textcolor{red}{50.6} & \textcolor{red}{47.6} & \textcolor{red}{35.2} & \textcolor{red}{33.8} & \textcolor{red}{32.0} & \textcolor{red}{34.1} & \textcolor{red}{33.7} \\
 & SFT & \textcolor{red}{65.8} & 66.9 & \textcolor{red}{65.8} & \textcolor{red}{65.9} & \textcolor{red}{66.2} & \textcolor{red}{62.0} & \textcolor{red}{60.4} & \textcolor{red}{65.6} & \textcolor{red}{62.3} & \textcolor{red}{64.6} & \textcolor{red}{57.8} & \textcolor{red}{58.2} & 62.0 & 61.6 & 62.2 \\
 & SFT+DPO & \textcolor{red}{64.0} & \textcolor{red}{65.8} & \textcolor{red}{66.6} & 68.0 & 68.9 & \textcolor{red}{62.0} & \textcolor{red}{64.4} & \textcolor{red}{64.8} & \textcolor{red}{63.2} & \textcolor{red}{63.5} & \textcolor{red}{57.8} & 63.0 & 63.0 & 61.8 & 65.0 \\
 & DPO & 67.3 & \textcolor{red}{65.6} & \textcolor{red}{65.4} & \textcolor{red}{65.0} & \textcolor{red}{65.3} & 68.2 & 69.1 & 71.5 & 70.3 & 72.2 & 66.1 & 66.8 & 69.7 & 68.6 & 69.0 \\
\midrule
\multirow{5}{*}{RoBERTa} & w/o synthetic data & \multicolumn{5}{c}{69.7} & \multicolumn{5}{c}{73.6} & \multicolumn{5}{c}{63.1} \\
 & Prompting & \textcolor{red}{66.5} & \textcolor{red}{59.3} & \textcolor{red}{58.0} & \textcolor{red}{55.5} & \textcolor{red}{54.6} & \textcolor{red}{59.0} & \textcolor{red}{57.3} & \textcolor{red}{57.1} & \textcolor{red}{55.2} & \textcolor{red}{54.9} & \textcolor{red}{49.1} & \textcolor{red}{42.2} & \textcolor{red}{40.2} & \textcolor{red}{38.3} & \textcolor{red}{36.8} \\
 & SFT & \textcolor{red}{69.3} & \textcolor{red}{69.5} & 70.1 & 70.3 & 70.1 & \textcolor{red}{66.9} & \textcolor{red}{69.4} & \textcolor{red}{69.4} & \textcolor{red}{68.6} & \textcolor{red}{68.2} & 65.5 & 66.8 & 66.6 & 66.8 & 66.7 \\
 & SFT+DPO & \textcolor{red}{67.9} & \textcolor{red}{68.8} & 70.3 & \textcolor{red}{69.0} & 70.6 & \textcolor{red}{68.8} & \textcolor{red}{69.9} & \textcolor{red}{70.5} & \textcolor{red}{68.1} & \textcolor{red}{68.6} & 65.1 & 67.6 & 67.7 & 66.1 & 66.7 \\
 & DPO & 70.7 & \textcolor{red}{69.5} & \textcolor{red}{68.7} & 70.0 & \textcolor{red}{68.1} & 77.1 & 74.9 & 76.9 & 75.8 & 74.7 & 69.2 & 69.9 & 71.5 & 72.0 & 71.6 \\
\midrule
\multirow{5}{*}{DeBERTa} & w/o synthetic data & \multicolumn{5}{c}{75.2} & \multicolumn{5}{c}{74.4} & \multicolumn{5}{c}{68.1} \\
 & Prompting & \textcolor{red}{66.2} & \textcolor{red}{58.2} & \textcolor{red}{53.7} & \textcolor{red}{52.5} & \textcolor{red}{49.9} & \textcolor{red}{56.1} & \textcolor{red}{54.8} & \textcolor{red}{53.6} & \textcolor{red}{53.5} & \textcolor{red}{53.2} & \textcolor{red}{43.6} & \textcolor{red}{38.4} & \textcolor{red}{41.4} & \textcolor{red}{39.8} & \textcolor{red}{37.7} \\
 & SFT & \textcolor{red}{70.8} & \textcolor{red}{70.3} & \textcolor{red}{70.2} & \textcolor{red}{70.7} & \textcolor{red}{69.4} & \textcolor{red}{68.9} & \textcolor{red}{66.7} & \textcolor{red}{67.7} & \textcolor{red}{68.3} & \textcolor{red}{65.9} & \textcolor{red}{64.4} & \textcolor{red}{63.6} & \textcolor{red}{64.8} & \textcolor{red}{63.0} & \textcolor{red}{63.0} \\
 & SFT+DPO & \textcolor{red}{71.4} & \textcolor{red}{70.4} & \textcolor{red}{69.8} & \textcolor{red}{69.8} & \textcolor{red}{72.5} & \textcolor{red}{70.5} & \textcolor{red}{67.2} & \textcolor{red}{69.2} & \textcolor{red}{69.5} & \textcolor{red}{66.3} & \textcolor{red}{66.7} & \textcolor{red}{65.4} & \textcolor{red}{67.9} & \textcolor{red}{66.3} & \textcolor{red}{65.7} \\
 & DPO & \textcolor{red}{72.8} & \textcolor{red}{68.9} & \textcolor{red}{71.7} & \textcolor{red}{70.8} & \textcolor{red}{71.9} & \textcolor{red}{74.4} & 74.9 & 76.0 & 75.8 & 74.6 & 70.5 & 71.5 & 73.0 & 72.4 & 72.2 \\
\bottomrule
\end{tabular}
\caption{F1 for correct answers ($\times$100) of each scorer trained on real data plus a random sample of synthetic answers from each synthesizer, by amount of synthetic data (1$\times$--5$\times$). Synthetic cells are means over five sampling draws; the per-level score without synthetic data spans each test set; \textcolor{red}{red} marks cells below it.}
\label{tab:appendix-random-amount-level-2}
\end{table*}

\begin{table*}[t]
\centering
\footnotesize
\setlength{\tabcolsep}{2pt}
\begin{tabular}{llccccccccccccccc}
\toprule
 & & \multicolumn{5}{c}{BEETLE UQ} & \multicolumn{5}{c}{SciEntsBank UQ} & \multicolumn{5}{c}{SciEntsBank UD} \\
\cmidrule(lr){3-7}\cmidrule(lr){8-12}\cmidrule(lr){13-17}
Model & Selection & 1$\times$ & 2$\times$ & 3$\times$ & 4$\times$ & 5$\times$ & 1$\times$ & 2$\times$ & 3$\times$ & 4$\times$ & 5$\times$ & 1$\times$ & 2$\times$ & 3$\times$ & 4$\times$ & 5$\times$ \\
\midrule
\multirow{9}{*}{ModernBERT} & w/o synthetic data & \multicolumn{5}{c}{68.6} & \multicolumn{5}{c}{64.6} & \multicolumn{5}{c}{64.2} \\
 & Random & 71.3 & 73.6 & 74.5 & 73.3 & 74.7 & 67.0 & 67.7 & 68.8 & 72.1 & 71.5 & 68.8 & 68.8 & 69.3 & 70.7 & 70.8 \\
 & Highest-Rank & 71.8 & -- & -- & -- & -- & 64.6 & -- & -- & -- & -- & \textcolor{red}{62.1} & -- & -- & -- & -- \\
 & Lowest-Rank & \textcolor{red}{59.9} & -- & -- & -- & -- & \textcolor{red}{47.1} & -- & -- & -- & -- & \textcolor{red}{45.0} & -- & -- & -- & -- \\
 & Hard Selection & 73.8 & 73.8 & 74.9 & 74.9 & 74.7 & 69.8 & 70.5 & 69.8 & 70.8 & 71.6 & 69.7 & 70.2 & 70.6 & 70.5 & 70.9 \\
 & $(0, 30)$ & 71.7 & 74.7 & 74.4 & 75.5 & 75.2 & 67.7 & 67.8 & 71.9 & 71.7 & 71.3 & 68.8 & 69.3 & 71.1 & 71.3 & 71.6 \\
 & $(0, 50)$ & 71.8 & 74.0 & 74.3 & 75.1 & 74.4 & 68.5 & 68.1 & 66.8 & 70.8 & 69.2 & 70.0 & 69.0 & 69.5 & 71.4 & 70.7 \\
 & $(10, 50)$ & 72.9 & 73.6 & 74.3 & 74.7 & 75.8 & 69.1 & 68.4 & 68.5 & 70.5 & 69.4 & 69.8 & 70.6 & 71.2 & 71.1 & 71.1 \\
 & $(20, 20)$ & 71.9 & 74.3 & 75.5 & 75.0 & 75.7 & 67.6 & 69.5 & 71.2 & 72.2 & 71.3 & 67.6 & 70.6 & 70.7 & 70.7 & 72.1 \\
\midrule
\multirow{9}{*}{BERT} & w/o synthetic data & \multicolumn{5}{c}{63.5} & \multicolumn{5}{c}{49.9} & \multicolumn{5}{c}{57.2} \\
 & Random & 67.4 & 69.1 & 69.2 & 70.0 & 70.5 & 53.8 & 58.3 & 58.8 & 62.0 & 63.6 & 59.7 & 62.2 & 63.8 & 64.4 & 63.3 \\
 & Highest-Rank & 67.2 & -- & -- & -- & -- & 54.5 & -- & -- & -- & -- & \textcolor{red}{52.5} & -- & -- & -- & -- \\
 & Lowest-Rank & \textcolor{red}{55.3} & -- & -- & -- & -- & \textcolor{red}{42.4} & -- & -- & -- & -- & \textcolor{red}{41.6} & -- & -- & -- & -- \\
 & Hard Selection & 69.8 & 71.2 & 72.3 & 72.6 & 72.6 & 55.3 & 59.7 & 62.4 & 63.3 & 63.2 & 62.2 & 65.4 & 66.3 & 68.0 & 67.0 \\
 & $(0, 30)$ & 70.3 & 70.5 & 71.3 & 71.9 & 71.8 & 55.2 & 59.4 & 61.1 & 62.5 & 63.8 & 61.6 & 64.9 & 66.3 & 67.4 & 68.2 \\
 & $(0, 50)$ & 68.4 & 70.9 & 71.2 & 72.2 & 72.8 & 54.8 & 57.1 & 60.9 & 61.0 & 60.6 & 60.0 & 63.9 & 65.6 & 66.7 & 66.0 \\
 & $(10, 50)$ & 69.3 & 71.2 & 71.7 & 72.0 & 72.9 & 53.7 & 58.6 & 61.2 & 60.8 & 61.3 & 61.2 & 66.2 & 66.4 & 66.9 & 67.1 \\
 & $(20, 20)$ & 69.0 & 69.1 & 72.0 & 71.6 & 72.2 & 55.9 & 59.3 & 60.6 & 61.9 & 63.4 & 60.3 & 63.4 & 65.9 & 66.9 & 66.2 \\
\midrule
\multirow{9}{*}{RoBERTa} & w/o synthetic data & \multicolumn{5}{c}{71.1} & \multicolumn{5}{c}{66.5} & \multicolumn{5}{c}{66.9} \\
 & Random & 73.9 & 75.5 & 74.7 & 75.9 & 74.4 & 68.4 & 69.5 & 68.2 & 69.6 & 67.5 & 68.1 & 69.3 & 69.2 & 69.4 & 68.9 \\
 & Highest-Rank & 71.5 & -- & -- & -- & -- & \textcolor{red}{61.7} & -- & -- & -- & -- & \textcolor{red}{58.7} & -- & -- & -- & -- \\
 & Lowest-Rank & \textcolor{red}{59.3} & -- & -- & -- & -- & \textcolor{red}{55.4} & -- & -- & -- & -- & \textcolor{red}{52.4} & -- & -- & -- & -- \\
 & Hard Selection & 72.4 & 73.4 & 73.9 & 74.5 & 74.6 & 68.3 & 69.5 & 70.2 & 70.9 & 71.6 & 69.1 & 71.1 & 70.8 & 70.7 & 70.4 \\
 & $(0, 30)$ & 71.8 & 74.4 & 74.8 & 75.4 & 75.2 & 69.7 & 70.7 & 69.8 & 71.3 & 72.1 & 70.3 & 71.1 & 71.3 & 71.9 & 71.6 \\
 & $(0, 50)$ & 72.0 & 72.4 & 74.2 & 75.6 & 75.0 & 67.2 & 68.1 & 68.2 & 69.6 & 68.7 & 68.6 & 70.4 & 70.8 & 71.4 & 70.6 \\
 & $(10, 50)$ & 72.0 & 72.8 & 75.2 & 75.0 & 75.7 & 67.4 & 68.3 & 67.6 & 69.1 & 67.6 & 69.4 & 70.8 & 70.3 & 71.1 & 71.3 \\
 & $(20, 20)$ & 71.5 & 74.6 & 75.6 & 75.3 & 76.1 & 68.8 & 70.3 & 71.0 & 71.8 & 70.1 & 68.9 & 71.3 & 70.8 & 71.1 & 69.1 \\
\midrule
\multirow{9}{*}{DeBERTa} & w/o synthetic data & \multicolumn{5}{c}{74.4} & \multicolumn{5}{c}{69.4} & \multicolumn{5}{c}{70.3} \\
 & Random & 75.2 & 75.4 & 74.9 & 75.6 & 76.3 & 69.9 & 71.5 & 71.1 & 72.7 & 71.9 & 70.5 & 70.7 & 71.9 & 71.8 & 71.5 \\
 & Highest-Rank & \textcolor{red}{72.1} & -- & -- & -- & -- & \textcolor{red}{67.6} & -- & -- & -- & -- & \textcolor{red}{66.5} & -- & -- & -- & -- \\
 & Lowest-Rank & \textcolor{red}{61.9} & -- & -- & -- & -- & \textcolor{red}{56.4} & -- & -- & -- & -- & \textcolor{red}{52.6} & -- & -- & -- & -- \\
 & Hard Selection & \textcolor{red}{74.1} & 74.9 & 75.2 & 75.3 & 76.1 & 69.5 & 71.5 & 71.2 & 72.0 & 71.7 & 71.1 & 72.8 & 73.4 & 73.5 & 71.7 \\
 & $(0, 30)$ & \textcolor{red}{73.6} & 75.3 & 76.0 & 75.5 & 76.4 & 70.0 & 71.4 & 70.8 & 72.9 & 71.8 & 71.2 & 73.0 & 72.7 & 73.5 & 72.5 \\
 & $(0, 50)$ & \textcolor{red}{71.7} & 75.6 & 75.1 & 76.6 & 77.1 & 70.0 & 70.8 & \textcolor{red}{69.4} & 71.8 & 69.9 & 72.9 & 72.6 & 72.2 & 72.1 & 72.0 \\
 & $(10, 50)$ & \textcolor{red}{73.3} & 75.0 & 76.4 & 76.7 & 77.0 & \textcolor{red}{69.0} & 69.8 & 70.2 & 70.0 & 70.4 & 72.7 & 72.6 & 72.0 & 72.5 & 72.0 \\
 & $(20, 20)$ & \textcolor{red}{74.1} & 76.2 & 76.1 & 76.9 & 77.6 & \textcolor{red}{68.8} & 72.1 & 73.4 & 73.5 & 72.2 & 73.5 & 73.3 & 72.5 & 73.0 & 72.1 \\
\bottomrule
\end{tabular}
\caption{F1 for incorrect answers ($\times$100) of each scorer with each selection policy, by amount of synthetic data (1$\times$--5$\times$). All policies select from DPO answers; cells are means over five sampling draws; Highest-Rank and Lowest-Rank were run at 1$\times$ only. The per-level score without synthetic data spans each test set; \textcolor{red}{red} marks cells below it.}
\label{tab:appendix-selection-amount-level-0}
\end{table*}

\begin{table*}[t]
\centering
\footnotesize
\setlength{\tabcolsep}{2pt}
\begin{tabular}{llccccccccccccccc}
\toprule
 & & \multicolumn{5}{c}{BEETLE UQ} & \multicolumn{5}{c}{SciEntsBank UQ} & \multicolumn{5}{c}{SciEntsBank UD} \\
\cmidrule(lr){3-7}\cmidrule(lr){8-12}\cmidrule(lr){13-17}
Model & Selection & 1$\times$ & 2$\times$ & 3$\times$ & 4$\times$ & 5$\times$ & 1$\times$ & 2$\times$ & 3$\times$ & 4$\times$ & 5$\times$ & 1$\times$ & 2$\times$ & 3$\times$ & 4$\times$ & 5$\times$ \\
\midrule
\multirow{9}{*}{ModernBERT} & w/o synthetic data & \multicolumn{5}{c}{47.8} & \multicolumn{5}{c}{42.0} & \multicolumn{5}{c}{34.4} \\
 & Random & 49.6 & 50.0 & 50.1 & 49.8 & 51.1 & 48.3 & 44.8 & 52.8 & 51.4 & 51.8 & 45.5 & 44.2 & 44.5 & 45.3 & 46.3 \\
 & Highest-Rank & 49.1 & -- & -- & -- & -- & 51.4 & -- & -- & -- & -- & 45.8 & -- & -- & -- & -- \\
 & Lowest-Rank & \textcolor{red}{32.9} & -- & -- & -- & -- & \textcolor{red}{34.1} & -- & -- & -- & -- & \textcolor{red}{27.0} & -- & -- & -- & -- \\
 & Hard Selection & 53.1 & 50.7 & 52.2 & 50.6 & 52.2 & 52.9 & 53.5 & 53.8 & 53.3 & 53.4 & 47.6 & 49.9 & 48.1 & 48.4 & 48.5 \\
 & $(0, 30)$ & 50.0 & 53.2 & 53.2 & 52.4 & 54.1 & 51.3 & 50.2 & 54.7 & 53.4 & 53.6 & 47.6 & 49.2 & 50.2 & 49.4 & 51.0 \\
 & $(0, 50)$ & 52.8 & 52.6 & 52.4 & 52.5 & 52.7 & 51.7 & 53.7 & 53.0 & 53.4 & 54.5 & 50.0 & 50.2 & 49.6 & 51.3 & 51.1 \\
 & $(10, 50)$ & 52.9 & 52.6 & 51.3 & 53.0 & 52.4 & 53.1 & 55.2 & 54.9 & 55.3 & 54.8 & 49.0 & 50.4 & 50.9 & 50.7 & 51.2 \\
 & $(20, 20)$ & 52.0 & 52.3 & 51.8 & 52.2 & 52.0 & 48.9 & 51.4 & 52.1 & 51.3 & 52.5 & 44.1 & 48.6 & 49.3 & 49.6 & 47.8 \\
\midrule
\multirow{9}{*}{BERT} & w/o synthetic data & \multicolumn{5}{c}{37.8} & \multicolumn{5}{c}{16.0} & \multicolumn{5}{c}{24.9} \\
 & Random & 49.9 & 50.9 & 51.0 & 50.5 & 49.9 & 47.1 & 46.3 & 46.6 & 47.2 & 46.9 & 40.5 & 42.6 & 40.7 & 42.2 & 40.6 \\
 & Highest-Rank & 49.2 & -- & -- & -- & -- & 46.0 & -- & -- & -- & -- & 42.1 & -- & -- & -- & -- \\
 & Lowest-Rank & \textcolor{red}{34.4} & -- & -- & -- & -- & 33.7 & -- & -- & -- & -- & 26.3 & -- & -- & -- & -- \\
 & Hard Selection & 52.5 & 52.7 & 52.7 & 53.6 & 50.8 & 46.8 & 48.5 & 49.2 & 49.8 & 50.0 & 44.3 & 44.7 & 44.9 & 46.3 & 45.6 \\
 & $(0, 30)$ & 53.7 & 54.7 & 55.6 & 54.3 & 53.9 & 45.5 & 47.4 & 47.2 & 49.8 & 50.3 & 42.2 & 46.6 & 46.9 & 47.0 & 47.1 \\
 & $(0, 50)$ & 53.1 & 54.6 & 54.8 & 53.9 & 52.8 & 49.4 & 49.6 & 51.7 & 49.8 & 51.6 & 44.0 & 47.0 & 48.7 & 49.0 & 49.5 \\
 & $(10, 50)$ & 54.0 & 54.4 & 55.1 & 53.4 & 53.5 & 48.6 & 49.0 & 50.8 & 49.6 & 51.7 & 45.1 & 47.7 & 48.3 & 47.9 & 48.3 \\
 & $(20, 20)$ & 52.2 & 53.2 & 55.0 & 53.7 & 55.3 & 43.5 & 46.2 & 45.3 & 45.9 & 49.9 & 40.9 & 43.2 & 43.5 & 43.9 & 44.2 \\
\midrule
\multirow{9}{*}{RoBERTa} & w/o synthetic data & \multicolumn{5}{c}{44.8} & \multicolumn{5}{c}{36.7} & \multicolumn{5}{c}{32.1} \\
 & Random & 51.7 & 51.3 & 49.9 & 51.9 & 50.8 & 52.3 & 49.2 & 52.4 & 50.5 & 49.4 & 42.6 & 42.5 & 43.7 & 42.3 & 40.2 \\
 & Highest-Rank & 50.7 & -- & -- & -- & -- & 48.1 & -- & -- & -- & -- & 43.5 & -- & -- & -- & -- \\
 & Lowest-Rank & \textcolor{red}{34.0} & -- & -- & -- & -- & \textcolor{red}{27.7} & -- & -- & -- & -- & \textcolor{red}{24.0} & -- & -- & -- & -- \\
 & Hard Selection & 51.0 & 52.4 & 50.6 & 51.3 & 50.6 & 50.5 & 51.1 & 52.3 & 52.8 & 53.9 & 44.5 & 45.4 & 45.1 & 44.8 & 45.3 \\
 & $(0, 30)$ & 53.8 & 53.8 & 53.4 & 54.2 & 54.1 & 52.3 & 51.8 & 53.7 & 54.5 & 55.1 & 45.1 & 45.9 & 46.7 & 46.4 & 45.6 \\
 & $(0, 50)$ & 52.8 & 53.0 & 53.5 & 53.1 & 52.8 & 51.6 & 51.3 & 54.1 & 52.9 & 54.5 & 47.1 & 47.1 & 48.6 & 48.5 & 48.2 \\
 & $(10, 50)$ & 52.8 & 53.3 & 53.7 & 52.9 & 52.7 & 48.7 & 51.2 & 53.9 & 53.5 & 54.5 & 46.6 & 47.4 & 47.3 & 48.2 & 47.8 \\
 & $(20, 20)$ & 50.7 & 54.4 & 53.6 & 53.6 & 51.5 & 49.7 & 51.3 & 52.5 & 53.0 & 54.7 & 44.8 & 45.6 & 45.4 & 45.0 & 44.4 \\
\midrule
\multirow{9}{*}{DeBERTa} & w/o synthetic data & \multicolumn{5}{c}{45.6} & \multicolumn{5}{c}{45.1} & \multicolumn{5}{c}{38.1} \\
 & Random & 51.1 & 49.9 & 51.0 & 50.8 & 51.2 & 52.4 & 53.3 & 52.1 & 53.0 & 54.0 & 47.0 & 45.0 & 44.1 & 44.0 & 44.1 \\
 & Highest-Rank & 48.7 & -- & -- & -- & -- & 53.4 & -- & -- & -- & -- & 46.3 & -- & -- & -- & -- \\
 & Lowest-Rank & \textcolor{red}{29.0} & -- & -- & -- & -- & \textcolor{red}{33.4} & -- & -- & -- & -- & \textcolor{red}{25.3} & -- & -- & -- & -- \\
 & Hard Selection & 52.3 & 51.0 & 51.6 & 50.8 & 51.1 & 53.0 & 54.0 & 55.0 & 52.5 & 56.1 & 47.1 & 48.7 & 47.8 & 47.5 & 45.6 \\
 & $(0, 30)$ & 52.4 & 52.4 & 53.5 & 53.5 & 53.2 & 52.6 & 52.7 & 56.3 & 54.6 & 56.3 & 49.0 & 50.6 & 49.4 & 48.6 & 47.9 \\
 & $(0, 50)$ & 51.8 & 52.5 & 52.1 & 51.5 & 51.0 & 53.5 & 53.9 & 55.2 & 56.2 & 55.9 & 48.6 & 49.8 & 49.9 & 49.9 & 50.7 \\
 & $(10, 50)$ & 52.4 & 51.5 & 52.5 & 51.7 & 51.0 & 54.2 & 53.9 & 56.7 & 55.3 & 57.6 & 48.7 & 48.6 & 49.2 & 49.8 & 49.9 \\
 & $(20, 20)$ & 51.2 & 52.2 & 52.8 & 53.5 & 52.2 & 51.1 & 52.4 & 54.2 & 54.1 & 55.7 & 49.0 & 47.6 & 46.7 & 46.0 & 45.6 \\
\bottomrule
\end{tabular}
\caption{F1 for partially correct answers ($\times$100) of each scorer with each selection policy, by amount of synthetic data (1$\times$--5$\times$). All policies select from DPO answers; cells are means over five sampling draws; Highest-Rank and Lowest-Rank were run at 1$\times$ only. The per-level score without synthetic data spans each test set; \textcolor{red}{red} marks cells below it.}
\label{tab:appendix-selection-amount-level-1}
\end{table*}

\begin{table*}[t]
\centering
\footnotesize
\setlength{\tabcolsep}{2pt}
\begin{tabular}{llccccccccccccccc}
\toprule
 & & \multicolumn{5}{c}{BEETLE UQ} & \multicolumn{5}{c}{SciEntsBank UQ} & \multicolumn{5}{c}{SciEntsBank UD} \\
\cmidrule(lr){3-7}\cmidrule(lr){8-12}\cmidrule(lr){13-17}
Model & Selection & 1$\times$ & 2$\times$ & 3$\times$ & 4$\times$ & 5$\times$ & 1$\times$ & 2$\times$ & 3$\times$ & 4$\times$ & 5$\times$ & 1$\times$ & 2$\times$ & 3$\times$ & 4$\times$ & 5$\times$ \\
\midrule
\multirow{9}{*}{ModernBERT} & w/o synthetic data & \multicolumn{5}{c}{69.0} & \multicolumn{5}{c}{73.8} & \multicolumn{5}{c}{65.3} \\
 & Random & \textcolor{red}{67.0} & \textcolor{red}{67.4} & \textcolor{red}{68.4} & \textcolor{red}{67.4} & \textcolor{red}{68.0} & \textcolor{red}{72.6} & \textcolor{red}{72.4} & 76.3 & \textcolor{red}{73.6} & 75.9 & 66.4 & 68.4 & 70.2 & 69.8 & 71.0 \\
 & Highest-Rank & \textcolor{red}{67.7} & -- & -- & -- & -- & \textcolor{red}{70.6} & -- & -- & -- & -- & \textcolor{red}{63.2} & -- & -- & -- & -- \\
 & Lowest-Rank & \textcolor{red}{50.2} & -- & -- & -- & -- & \textcolor{red}{39.2} & -- & -- & -- & -- & \textcolor{red}{44.1} & -- & -- & -- & -- \\
 & Hard Selection & 71.8 & 69.7 & 71.0 & 69.7 & 69.6 & 75.4 & 76.5 & 75.7 & 75.2 & 74.2 & 69.2 & 70.6 & 70.3 & 69.7 & 70.3 \\
 & $(0, 30)$ & \textcolor{red}{68.5} & 71.0 & 71.1 & 70.1 & 70.0 & 74.5 & \textcolor{red}{72.4} & 75.8 & 74.8 & 74.9 & 68.0 & 69.5 & 70.3 & 69.9 & 70.0 \\
 & $(0, 50)$ & 70.3 & \textcolor{red}{68.8} & 69.1 & 69.4 & \textcolor{red}{68.5} & 74.1 & 75.1 & 75.8 & \textcolor{red}{73.4} & 75.9 & 67.9 & 68.7 & 70.6 & 68.5 & 70.5 \\
 & $(10, 50)$ & 69.7 & 69.7 & \textcolor{red}{68.8} & 69.4 & \textcolor{red}{69.0} & 75.0 & 74.8 & 76.2 & 75.4 & 76.6 & 69.3 & 69.4 & 69.3 & 69.6 & 70.2 \\
 & $(20, 20)$ & 69.4 & 71.8 & 70.2 & 70.1 & 70.6 & \textcolor{red}{72.6} & 77.3 & 75.5 & 75.4 & 76.8 & 67.7 & 71.8 & 72.0 & 69.9 & 71.8 \\
\midrule
\multirow{9}{*}{BERT} & w/o synthetic data & \multicolumn{5}{c}{66.7} & \multicolumn{5}{c}{65.7} & \multicolumn{5}{c}{60.6} \\
 & Random & 67.3 & \textcolor{red}{65.6} & \textcolor{red}{65.4} & \textcolor{red}{65.0} & \textcolor{red}{65.3} & 68.2 & 69.1 & 71.5 & 70.3 & 72.2 & 66.1 & 66.8 & 69.7 & 68.6 & 69.0 \\
 & Highest-Rank & \textcolor{red}{64.3} & -- & -- & -- & -- & 66.3 & -- & -- & -- & -- & 62.8 & -- & -- & -- & -- \\
 & Lowest-Rank & \textcolor{red}{39.3} & -- & -- & -- & -- & \textcolor{red}{49.1} & -- & -- & -- & -- & \textcolor{red}{42.5} & -- & -- & -- & -- \\
 & Hard Selection & 69.7 & 69.8 & 69.5 & 69.2 & 69.5 & 68.6 & 71.4 & 71.7 & 72.8 & 73.3 & 65.8 & 68.7 & 69.4 & 69.8 & 70.3 \\
 & $(0, 30)$ & 68.5 & 69.5 & 69.8 & 69.2 & 68.9 & 68.2 & 68.5 & 71.4 & 73.0 & 72.6 & 67.6 & 70.0 & 69.6 & 69.7 & 70.4 \\
 & $(0, 50)$ & 68.3 & 68.3 & 68.3 & 67.4 & 67.2 & 67.8 & 70.8 & 72.5 & 72.5 & 72.6 & 66.7 & 67.6 & 69.8 & 69.9 & 70.4 \\
 & $(10, 50)$ & 68.9 & 69.6 & 68.0 & 69.0 & 68.1 & 68.1 & 71.2 & 72.0 & 72.2 & 72.8 & 66.3 & 68.9 & 70.5 & 70.1 & 70.1 \\
 & $(20, 20)$ & 67.3 & 67.8 & 70.7 & 68.9 & 70.3 & 67.8 & 71.1 & 70.2 & 71.9 & 73.0 & 66.9 & 70.2 & 70.5 & 70.8 & 71.9 \\
\midrule
\multirow{9}{*}{RoBERTa} & w/o synthetic data & \multicolumn{5}{c}{69.7} & \multicolumn{5}{c}{73.6} & \multicolumn{5}{c}{63.1} \\
 & Random & 70.7 & \textcolor{red}{69.5} & \textcolor{red}{68.7} & 70.0 & \textcolor{red}{68.1} & 77.1 & 74.9 & 76.9 & 75.8 & 74.7 & 69.2 & 69.9 & 71.5 & 72.0 & 71.6 \\
 & Highest-Rank & \textcolor{red}{68.7} & -- & -- & -- & -- & \textcolor{red}{71.9} & -- & -- & -- & -- & \textcolor{red}{62.1} & -- & -- & -- & -- \\
 & Lowest-Rank & \textcolor{red}{48.4} & -- & -- & -- & -- & \textcolor{red}{58.0} & -- & -- & -- & -- & \textcolor{red}{53.6} & -- & -- & -- & -- \\
 & Hard Selection & 71.2 & \textcolor{red}{69.1} & 70.3 & 70.3 & 70.1 & 75.6 & 77.3 & 76.1 & 76.5 & 75.7 & 69.8 & 69.5 & 70.9 & 70.9 & 71.0 \\
 & $(0, 30)$ & 71.6 & 71.1 & 70.9 & 70.9 & 71.1 & 76.2 & 75.7 & 77.3 & 77.9 & 77.2 & 70.7 & 70.4 & 71.6 & 71.1 & 71.6 \\
 & $(0, 50)$ & \textcolor{red}{69.3} & \textcolor{red}{69.7} & \textcolor{red}{69.7} & 70.2 & \textcolor{red}{69.5} & 74.6 & 75.8 & 76.2 & 75.6 & 77.1 & 69.4 & 69.9 & 69.7 & 69.8 & 71.1 \\
 & $(10, 50)$ & \textcolor{red}{69.1} & \textcolor{red}{69.7} & 70.9 & 70.8 & \textcolor{red}{69.3} & 75.1 & 75.8 & 77.2 & 76.4 & 77.7 & 69.4 & 70.1 & 70.6 & 70.4 & 71.3 \\
 & $(20, 20)$ & 70.7 & 71.8 & 72.2 & 72.6 & 69.8 & 76.5 & 77.4 & 78.2 & 77.5 & 77.5 & 70.6 & 72.1 & 72.0 & 72.1 & 71.9 \\
\midrule
\multirow{9}{*}{DeBERTa} & w/o synthetic data & \multicolumn{5}{c}{75.2} & \multicolumn{5}{c}{74.4} & \multicolumn{5}{c}{68.1} \\
 & Random & \textcolor{red}{72.8} & \textcolor{red}{68.9} & \textcolor{red}{71.7} & \textcolor{red}{70.8} & \textcolor{red}{71.9} & \textcolor{red}{74.4} & 74.9 & 76.0 & 75.8 & 74.6 & 70.5 & 71.5 & 73.0 & 72.4 & 72.2 \\
 & Highest-Rank & \textcolor{red}{68.5} & -- & -- & -- & -- & \textcolor{red}{73.3} & -- & -- & -- & -- & \textcolor{red}{68.0} & -- & -- & -- & -- \\
 & Lowest-Rank & \textcolor{red}{50.4} & -- & -- & -- & -- & \textcolor{red}{51.4} & -- & -- & -- & -- & \textcolor{red}{49.8} & -- & -- & -- & -- \\
 & Hard Selection & \textcolor{red}{73.4} & \textcolor{red}{71.1} & \textcolor{red}{72.3} & \textcolor{red}{72.9} & \textcolor{red}{72.2} & 74.6 & 75.3 & 76.8 & 75.9 & 76.2 & 69.9 & 71.1 & 72.2 & 71.1 & 72.3 \\
 & $(0, 30)$ & \textcolor{red}{71.0} & \textcolor{red}{72.0} & \textcolor{red}{72.2} & \textcolor{red}{72.6} & \textcolor{red}{72.5} & \textcolor{red}{74.2} & \textcolor{red}{74.3} & 76.4 & 74.9 & 77.3 & 69.7 & 72.6 & 73.0 & 71.5 & 72.4 \\
 & $(0, 50)$ & \textcolor{red}{70.6} & \textcolor{red}{72.8} & \textcolor{red}{70.7} & \textcolor{red}{70.5} & \textcolor{red}{71.5} & 75.8 & 75.8 & 75.3 & 75.8 & 76.3 & 71.1 & 70.3 & 71.5 & 70.6 & 72.6 \\
 & $(10, 50)$ & \textcolor{red}{71.4} & \textcolor{red}{71.6} & \textcolor{red}{71.7} & \textcolor{red}{71.5} & \textcolor{red}{71.1} & 76.7 & 74.7 & 77.2 & 76.7 & 77.3 & 71.3 & 70.7 & 71.7 & 72.4 & 72.7 \\
 & $(20, 20)$ & \textcolor{red}{72.1} & \textcolor{red}{72.6} & \textcolor{red}{72.6} & \textcolor{red}{72.2} & \textcolor{red}{72.1} & 76.4 & 75.6 & 76.6 & 75.8 & 77.7 & 70.8 & 73.3 & 74.1 & 73.5 & 73.8 \\
\bottomrule
\end{tabular}
\caption{F1 for correct answers ($\times$100) of each scorer with each selection policy, by amount of synthetic data (1$\times$--5$\times$). All policies select from DPO answers; cells are means over five sampling draws; Highest-Rank and Lowest-Rank were run at 1$\times$ only. The per-level score without synthetic data spans each test set; \textcolor{red}{red} marks cells below it.}
\label{tab:appendix-selection-amount-level-2}
\end{table*}

